\documentclass[11pt]{article}
\usepackage[margin=2.5cm]{geometry}
\usepackage{amsmath, amssymb, amsfonts}
\usepackage[none]{hyphenat} %(c) To prevet LaTeX form using hyphenated words
\usepackage{fancyhdr} %(c) This allows me to create a custom header and footer
\usepackage{graphicx}
\usepackage{float}
\usepackage{color}
\usepackage[x11names]{xcolor}
% \usepackage[nottoc, notlot, notlof]{tocbibind} %(c) for \tableofcontents .

\pagestyle{fancy} 
%(c) To clear the default header and the default hooter
\fancyhead{}
\fancyfoot{}
\fancyhead[L]{\slshape \MakeUppercase{Place Title Here}} %(c) We used \MakeUppercase{} to make my title capitalized and \slshape{} to make it in italics.  
\fancyhead[R]{\slshape Muhammad Ali Qattan}
\fancyfoot[C]{\thepage}
% \renewcommand{\headrulewidth}{0pt} %(c) That removes the horozilntal line of the header
% \renewcommand{\footrulewidth}{0.5pt} %(c) That appears the horozilntal line of the footer

% \parindent 0px %(c) 0px because I don't want the paragraph indentation 
\setlength{\parindent}{4em} %(c) also you can use this commad to change the width of the indentation
\setlength{\parskip}{1em} %(c) This to change the spacing between your paragraphs
\renewcommand{\baselinestretch}{1} %(c) That affencting the lines spacing

\begin{document}
  \begin{titlepage}
    \begin{center}
      \vspace*{1cm}
      \Large{\textbf{IB Mathematics SL}}\\
      \Large{\textbf{Internal Assessment}}\\
      \vfill %(c) To add amount of space between these two items (before and after it) to fill the page 
      \line(1,0){400}\\[1mm]
      \huge{\textbf{This is a Sample Title}}\\[3mm]
      \Large{\textbf{- This is a Sample Subtitle - }}\\[1mm]
      \line(1,0){400}\\
      \vfill %(c) to center them
      By Muhammad Ali Qattan\\
      Candidate 207047 \\[5mm]
      \today\\
    \end{center}
  \end{titlepage}

\tableofcontents %(c) for the table of contents
\thispagestyle{empty} %(c) remove the header and the footer from this page 
\clearpage %(c) clear tha page from any thing excipt the table of contents 

\setcounter{page}{1}

\section{Introduction}
In this report, we explore several key aspects of \colorbox{yellow}{mathematical writing and research communication}. 
The goal is to practice structuring a document using \LaTeX{}, while also understanding the 
criteria typically used to assess \textcolor{red}{mathematical investigations}. 

Throughout this document, we will present explanations, examples, and short exercises that 
you can modify while learning. We also include tables, equations, and references to other 
sections. \footnote{An example footnote}

\section{Scoring Criteria}
This section summarizes the criteria commonly used to evaluate mathematical work. 
Each subsection describes a different scoring component.

    \subsection{Communication}
    (See this table \ref{tab:data1} if you want)\\ 
    Clear communication is essential in mathematical writing. A well-structured report should:
        \begin{itemize}
            \item present ideas in a logical sequence,
            \item define all mathematical symbols,
            \item include diagrams or tables when helpful,
            \item use consistent notation throughout.
        \end{itemize}

    As an example, consider the following table summarizing notation used in this document:

        \begin{table}[h!]
            \centering
                \begin{tabular}{|c|c|}
                \hline
                Symbol & Meaning \\ \hline
                $a_n$  & $n$-th term of a sequence \\ \hline
                $f(x)$ & Function depending on $x$ \\ \hline
                $\sum$ & Summation notation \\ \hline
                \end{tabular}
            \caption{Common Mathematical Symbols}
            \label{tab:data1} %(c) To refer to this table later (we refered to this table in the line 64), for figures you can use {fig:thing}.
        \end{table}


\subsection{Mathematical Presentation}
Your work should be mathematically accurate and logically reasoned. 
This includes showing all steps and using appropriate formatting for equations.

Below is an example formatted using the \texttt{equation} environment:
\begin{equation}
    \int_0^1 x^2 \, dx = \left[ \frac{x^3}{3} \right]_0^1 = \frac{1}{3}.
\end{equation}

We can also typeset systems of equations:
\[
\begin{cases}
x + y = 7, \\
2x - y = 3.
\end{cases}
\]

\subsection{Personal Engagement}
Personal engagement means showing your own thinking, discoveries, or motivation. 
For example:
\begin{quote}
    ``I became curious about how the Fibonacci sequence grows, so I decided 
    to model it using exponential approximations.''
\end{quote}

You can also include simple experiments you perform, like generating numbers with a small script or manually exploring patterns.

\subsection{Reflection}
Reflection means analyzing what you have learned, what challenges you faced, and how the mathematics connects to real life.

Example reflective paragraph:

Mathematical modeling helped me understand how quickly exponential functions grow. 
I initially expected linear-like behavior, but after plotting several values, I realized 
that exponential growth is far more rapid. This changed how I interpret graphs in real-world contexts such as population models.

\subsection{Use of Mathematics}
This criterion assesses the correctness and depth of your mathematical work. 

You can practice by adding any of the following:
    \begin{itemize}
        \item derivations,
        \item algebraic manipulation,
        \item graphs created with TikZ,
        \item proofs or justifications.
    \end{itemize}

Here is a small example involving sequences:
\[
a_n = 2n + 1 \quad \Rightarrow \quad S_n = \sum_{k=1}^{n} (2k+1) = n^2 + 2n.
\]

\subsection{Communication}
This section is repeated in your original structure, but you may keep it if needed 
for additional practice. Here, we include a short paragraph exercise:

Effective communication ensures the reader can follow your mathematical arguments without confusion. Using clear paragraph transitions, labeled figures, and meaningful examples 
helps achieve a professional writing style.

\section{Conclusion}
In this practice document, we created a structured \LaTeX{} file containing sections, 
subsections, tables, equations, and lists. By modifying the content, adding new examples, 
and experimenting with formatting, you will gradually become more confident using 
\LaTeX{} for academic or professional writing. \cite{DBHS1}

Future work could involve adding bibliographies, figures, TikZ diagrams, or exploring 
packages such as \texttt{amsmath}, \texttt{graphicx}, and \texttt{hyperref}.

\pagebreak %(c) Or you can use \newpage .
\begin{thebibliography}{} 
    \bibitem{DBHS1} %(c) DBHS1 is name you can choose any name you want for refering to it (like in line 152)
    Alcosser, Howard.
    ``Diamond Bar High School.''
    \textit{Internal assessment: Mathematical Exploration.}
    Web. 26/11/2025.
    \bibitem{DBHS2}
    Alcosser, Howard.
    ``Diamond Bar High School.''
    \textit{Internal assessment: Mathematical Exploration.}
    Web. 26/11/2025.
    \bibitem{DBHS3}
    Alcosser, Howard.
    ``Diamond Bar High School.''
    \textit{Internal assessment: Mathematical Exploration.}
    Web. 26/11/2025.
\end{thebibliography}



\end{document}